\chapter{Research Methodology}
\label{ch:methodology}

\section{Research Design}

The study is a retrospective, observational prediction-modelling study on
de-identified critical-care data. It follows a pre-registered protocol
(\texttt{00\_preregistration.md}, locked before any modelling) and reports
in conformity with TRIPOD+AI \parencite{collins2024tripodai} and PROBAST+AI
\parencite{moons2025probastai}. The central methodological innovation is
the elevation of label operationalisation, validation split design, and
evaluation metric selection from nuisance parameters to primary objects of
study.

All code is implemented in R (v4.x) and Python (v3.11), with a single
configuration module (\texttt{config.R}) as the sole source of truth for paths,
clinical item identifiers, antibiotic name lists, and Sepsis-3 thresholds.

%===========================================================================
\section{Data Sources}
\label{sec:data}
%===========================================================================

\subsection{Development Data: MIMIC-IV v3.1}

The primary development data is MIMIC-IV v3.1 \parencite{johnson2023mimiciv},
a large, de-identified critical-care database comprising records from Beth Israel
Deaconess Medical Center (BIDMC), Boston. It was accessed under the required
PhysioNet credentialing agreement. The full database spans 8 clinical tables
used in this study (Table~\ref{tab:data_sources}).

\begin{table}[H]
  \centering
  \begin{tabular}{lrr}
    \toprule
    \textbf{Table} & \textbf{Exists} & \textbf{Size (MB)} \\
    \midrule
    admissions       & yes & 19.9 \\
    patients         & yes & 2.8 \\
    icustays         & yes & 3.3 \\
    chartevents      & yes & 3{,}502.6 \\
    labevents        & yes & 2{,}593.0 \\
    prescriptions    & yes & 606.7 \\
    microbiologyevents & yes & 117.2 \\
    inputevents      & yes & 401.3 \\
    \bottomrule
  \end{tabular}
  \caption{MIMIC-IV v3.1 source tables used in the V2 pipeline.}
  \label{tab:data_sources}
\end{table}

\subsection{External Validation Data: eICU-CRD v2.0}

External validation uses the eICU Collaborative Research Database v2.0
\parencite{pollard2018eicu}, a multi-centre collection spanning more than 200
US ICUs. eICU-CRD is independently collected, uses different electronic health
record systems, and represents a broader range of clinical practices and patient
populations than the single-centre MIMIC-IV, making it a meaningful external test.

%===========================================================================
\section{Cohort Construction and Inclusion Criteria}
%===========================================================================

The base cohort is constructed by applying the following inclusion criteria
uniformly across all three label variants:

\begin{itemize}
  \item Age $\geq 18$ years at ICU admission.
  \item First ICU stay per patient only (to avoid within-patient correlation from
  repeated admissions).
  \item ICU length of stay $\geq 4$ hours (to ensure sufficient data for
  hourly feature construction and at least one prediction opportunity).
\end{itemize}

The resulting base cohort contains \textbf{65{,}078 stays} in MIMIC-IV,
covering all stays identifiable across the three label variants.

%===========================================================================
\section{Outcome Definition: Three Pre-registered Sepsis-3 Label Variants}
\label{sec:label_variants}
%===========================================================================

Following Cohen et al.\ \parencite{cohen2024subtle} and Johnson et al.\
\parencite{johnson2018comparative}, the Sepsis-3 onset label is operationalised
as three independent but plausible readings of the consensus definition
\parencite{singer2016sepsis3}. The variants are pre-specified and locked before
any model fitting; they differ in the antibiotic-culture time windows and in the
SOFA baseline reference (Table~\ref{tab:label_variants}).

\begin{table}[H]
  \centering
  \small
  \begin{tabularx}{\textwidth}{lXrrrrl}
    \toprule
    \textbf{ID} & \textbf{Name} & \textbf{ABX$\to$cx (h)} & \textbf{cx$\to$ABX (h)} &
      \textbf{SOFA$-$ (h)} & \textbf{SOFA$+$ (h)} & \textbf{Baseline} \\
    \midrule
    A & Narrow          & 24 & 24 & 24 & 12 & Admission SOFA \\
    B & Seymour-Standard & 72 & 24 & 48 & 24 & Admission SOFA \\
    C & Liberal         & 72 & 72 & 48 & 24 & Rolling 24\,h min SOFA \\
    \bottomrule
  \end{tabularx}
  \caption{Three pre-registered Sepsis-3 label variants. ABX = antibiotic;
  cx = culture; SOFA$-$/SOFA$+$ = hours before/after suspicion time in which
  an acute SOFA rise of $\geq 2$ points must be observed.
  Variant B matches the original Seymour et al.\ (2016) window specification
  and is designated the primary analysis.}
  \label{tab:label_variants}
\end{table}

All three variants operationalise suspected infection as a qualifying antibiotic
order paired with a body-fluid culture within the specified windows. Surveillance
cultures (e.g.\ MRSA screens, CRE screens) are excluded. SOFA organ-dysfunction
components are computed from chartevents, labevents, inputevents, and
outputevents using the standard MIMIC-IV item identifiers.

Variant B (Seymour-Standard) is designated the \textbf{primary analysis}; Variants
A and C serve as pre-specified sensitivity analyses for the variance decomposition.

%===========================================================================
\section{Prediction Framing}
\label{sec:framing}
%===========================================================================

Following Lauritsen et al.\ \parencite{lauritsen2021framing} and Reyna et al.\
\parencite{reyna2020utility}:

\begin{itemize}
  \item \textbf{Observation window:} expanding from ICU admission (all data up
  to the current hour, forward-filled within each stay).
  \item \textbf{Prediction horizon:} 6 hours (will the patient develop sepsis
  onset within the next 6 hours?).
  \item \textbf{Window shift:} hourly.
  \item \textbf{Prediction trigger:} every at-risk person-hour from hour 1 of
  ICU admission.
  \item \textbf{Minimum observation:} 1 hour post-admission before eligibility.
  \item \textbf{Right-censoring:} 72 hours post-ICU-admission.
\end{itemize}

This framing matches the PhysioNet/CinC 2019 Challenge \parencite{reyna2020utility}
and the Lauritsen et al.\ design.

%===========================================================================
\section{Person-Hour Table: Competing Events}
\label{sec:person_hours}
%===========================================================================

Each cohort stay is expanded into a dense (stay$\_$id, hour) table with one row
per person-hour at risk. The competing-event structure follows Heyard et al.\
\parencite{heyard2019dynamic}:

\begin{itemize}
  \item event = 0: in-progress (censored at this hour)
  \item event = 1: sepsis onset in this hour (cause of interest)
  \item event = 2: discharge alive without sepsis (competing event)
  \item event = 3: death without sepsis (competing event)
\end{itemize}

Treating discharge and death as ordinary censoring overstates onset incidence
and biases parameter estimation. The multinomial cause-specific formulation
eliminates this bias by construction.

Treatment-anchor timestamps (first antibiotic order, first culture draw) are
extracted per stay to support treatment-anchored evaluation
\parencite{kamran2024evaluation}.

%===========================================================================
\section{Feature Construction}
\label{sec:features}
%===========================================================================

Features are extracted at each person-hour from the chartevents, labevents,
inputevents, and outputevents tables using the item identifiers in
\texttt{config.R}. The feature set comprises:

\begin{enumerate}
  \item \textbf{Vital signs} (last-observation-carried-forward per hour):
  heart rate, respiratory rate, SpO$_2$, MAP (invasive/noninvasive), temperature
  ($^\circ$C), GCS (eye + verbal + motor sum).
  \item \textbf{Laboratory values} (LOCF): lactate, creatinine, total bilirubin,
  platelet count, WBC, P/F ratio.
  \item \textbf{Vasopressor indicators:} any vasopressor (binary); any
  norepinephrine or epinephrine (binary).
  \item \textbf{Six-hour rolling slopes:} respiratory rate, MAP, lactate.
  \item \textbf{Missingness indicators:} one binary flag per laboratory value
  indicating whether the test was ordered in the current hour
  (lactate, creatinine, bilirubin, platelets, WBC, P/F ratio).
\end{enumerate}

Temperature recorded in Fahrenheit is converted to Celsius. FiO$_2$ values
$> 1.0$ (percent format) are divided by 100.

%===========================================================================
\section{Model Classes}
\label{sec:models}
%===========================================================================

\subsection{Primary: Dynamic Discrete-Time Competing-Risks Hazard Model}

Following D'Agostino et al.\ \parencite{dagostino1990pooled} and Heyard et al.\
\parencite{heyard2019dynamic}, the primary model is a multinomial pooled logistic
regression over person-hours:

\begin{equation}
  \lambda_r(t \mid \mathbf{X}_{it}) =
    P(T_i = t,\, R_i = r \mid T_i \geq t,\, \mathbf{X}_{it})
  = \frac{\exp\!\bigl(\alpha_r(t) + \boldsymbol{\beta}_r^\top \mathbf{X}_{it}\bigr)}
         {1 + \sum_{s} \exp\!\bigl(\alpha_s(t) + \boldsymbol{\beta}_s^\top \mathbf{X}_{it}\bigr)}
  \label{eq:primary_model}
\end{equation}

where $r \in \{1=\text{sepsis},\, 2=\text{discharge},\, 3=\text{death}\}$,
$\alpha_r(t)$ is the cause-specific baseline log-hazard modelled as a restricted
natural cubic spline in hours-since-admission with 5 knots
(at $t \in \{4, 13, 24, 38, 58\}$ hours, boundary knots at $\{0, 72\}$),
and $\boldsymbol{\beta}_r$ are cause-specific covariate coefficients.

Standard errors are cluster-robust sandwich estimates (R package
\texttt{sandwich}) clustered on \texttt{stay\_id} to account for
within-patient correlation across person-hours. The model is fitted via
\texttt{nnet::multinom()}.

\paragraph{Equivalence to Cox.} D'Agostino et al.\ prove that pooled logistic
regression over short equal-length intervals converges to time-dependent-covariate
Cox regression as the interval length approaches zero and the per-interval event
probability approaches zero. Both conditions are satisfied for one-hour ICU bins
with a per-hour onset probability below 0.5\%. The proposed model is therefore
the Cox model implemented in the form the data actually takes, without imposing
the proportional-hazards assumption.

\subsection{Mandatory ML Comparator: Gradient-Boosted Trees}

A gradient-boosted trees model (R package \texttt{xgboost}) is fitted under
identical conditions per label variant and temporal split. Maximum depth 6,
learning rate 0.1, early stopping on an internal validation set (15\% of
training rows), log-loss objective, column sampling by tree 0.8. Feature
importance is reported as gain.

\subsection{Static Cox Proportional-Hazards Model}

A standard Cox PH model is fitted using the first non-zero person-hour per stay
(static, baseline-covariate formulation) as a deliberately weak baseline to
quantify the cost of the field's conventional choice. Schoenfeld residual testing
is applied.

\subsection{Rule-Based Comparators}

NEWS2, qSOFA, and SIRS are computed at each person-hour from the same extracted
feature set. NEWS2 is the headline comparator following Evans et al.\
\parencite{evans2021ssc}.

%===========================================================================
\section{Validation Architecture}
\label{sec:validation}
%===========================================================================

\begin{table}[H]
  \centering
  \small
  \begin{tabularx}{\textwidth}{lXl}
    \toprule
    \textbf{Split} & \textbf{Design} & \textbf{Purpose} \\
    \midrule
    Internal (primary) & Temporal: \texttt{anchor\_year\_group}
      train = 2008--2016, test = 2017--2019
      \parencite{guo2022temporal} & Honest internal estimate \\
    Internal (secondary) & Random patient-level 75/25
      (seed = 42) & Quantify split optimism (H4) \\
    Treatment-anchored & Restricted to hours before first
      culture or antibiotic \parencite{kamran2024evaluation} & Remove
      clinician-suspicion leakage (H3) \\
    External (primary) & eICU-CRD v2.0
      \parencite{pollard2018eicu} & Generalisation (H5) \\
    \bottomrule
  \end{tabularx}
  \caption{Validation architecture. All splits are patient-level (no patient
  appears in both training and test sets). The primary analysis uses the
  temporal split.}
  \label{tab:validation}
\end{table}

%===========================================================================
\section{Evaluation Metrics}
\label{sec:metrics}
%===========================================================================

\subsection{Primary Metrics}

\textbf{PhysioNet Utility Score} \parencite{reyna2020utility} is the primary
discriminative metric. The score rewards early true positives (alerts issued
within 6 hours before onset), penalises late true positives and false negatives,
and penalises false positives. The score is normalised so that perfect prediction
scores 1.0, an optimal constant prediction scores 0.0, and a model can score
negative. A negative utility score is categorically worse than issuing no alert
at all. AUROC is designated \textbf{descriptive only} following Wang et al.\
\parencite{wang2025srpm}.

\textbf{Calibration} is assessed via calibration intercept, calibration slope,
and flexible calibration curve. Cause-specific calibration for the competing-risks
model follows Heyard et al.\ \parencite{heyard2020evaluation}.

\subsection{Secondary Metrics}

AUPRC (mandatory at low per-hour prevalence), alert burden (false alerts per
true alert; benchmark: 1.4 from Moor et al.\ \parencite{moor2023review}),
lead time relative to clinical recognition (benchmark: 3.7\,h from Moor et al.),
and decision-curve net benefit \parencite{vickers2006dca}.

%===========================================================================
\section{Sample Size}
\label{sec:sample_size}
%===========================================================================

Formal sample size calculation follows Riley et al.\ \parencite{riley2019sample}
(\texttt{pmsampsize} R package), targeting shrinkage factor $\geq 0.9$ and
apparent-vs-adjusted Nagelkerke $R^2$ difference $\leq 0.05$. The anticipated
AUC is 0.846 from Moor et al.\ \parencite{moor2023review}. Twenty-five
candidate predictors are specified. Sufficiency is evaluated per label variant;
underpowered variants are a finding, not a failure.

%===========================================================================
\section{Pre-registration and Reporting Standards}
%===========================================================================

The full protocol was pre-registered as \texttt{00\_preregistration.md} before
any data extraction. All six hypotheses (H1--H6), the three label variants, the
model classes, the validation splits, the primary metrics, and the go/no-go gates
are specified in that document. Post-hoc analyses are labelled as such. The
pipeline is released under a permissive open-source licence following the
TRIPOD+AI open-pipeline recommendation \parencite{collins2024tripodai}.
